\begin{center}
\textsc{Gauss} to \textsc{Gerling}.\quad Göttingen, 21~February 1822.
\end{center}

\par\noindent\hspace*{\parindent}\mbox{.\ .\ .\ .\ .\ .}\quad If the connections to the Göttingen observatory and to the Inselsberg cannot be made from one spot on the Meisner, you can now use \textls[120]{both} spots.
It adds very little work, for according to the principles of probability theory, if you measure at both points the angles between the objects visible from both, you need only repeat the measurements half as often at each point as you would have done if you had used only one station.
You will have no difficulty determining the relative positions of the two points with the greatest precision and yet without needless \mbox{effort.\quad.\ .\ .\ .\ .\ .}
